\honest{What Do We Mean By ``Rationality''?}{Eliezer Yudkowsky}{2009}

I mean two things by rationality. ``Epistemic rationality'' is making your beliefs more accurate. ``Instrumental rationality'' is getting what you value, which I sometimes call ``winning.'' I give no source for the pair. \nb{It is standard in philosophy. The post does not ask the obvious question about it: what to do when a false belief would get you what you value.}

(Truth does not mean certainty. We can only make our beliefs more probably accurate.)

Then I raise an objection myself. People say ``X is rational!'' when they mean ``I think X is true'' or ``I think X is good.'' So why have the word at all? My answer is that we need it for general statements, and my example is that rational agents maximize expected utility.

Next, the psychologists. People rate ``Bill is an accountant who plays jazz'' as more likely than ``Bill plays jazz.'' I leave out the description of Bill that the subjects were given, and point back to earlier essays for it. \nb{That description is what made ``accountant'' tempting, so without it the subjects look stranger than they were.} The judgment breaks a law of probability. The ``gold standards'' are probability theory and decision theory, and someone who obeys them is ``acting like a Bayesian.'' That is not acting like Mr.\ Spock, which does not help ``one hair''; a footnote extends rationality to ``urges, hunches, perceptions, and wordless intuitions.''

Then two admissions. First, no one can actually calculate and obey the mathematics I have just named as the gold standard. So the gold standard is not something a person can follow. There is a further art, and I list its parts: learn your flaws, overcome your biases, get into good emotional shape, ``et cetera, et cetera.'' Second, the mathematics itself is disputed in some puzzles, such as Newcomb's problem.

It is futile to settle a dispute by defining ``rational'' so that your answer wins. If you say ``It's (epistemically) rational for me to believe X, but the truth is Y,'' you are probably using the word to mean something else. I do the same for action: if the rational thing and the right thing differ, you are almost certainly using one of the two words in some other sense.

I use the word ``normatively, to pick out desirable patterns of thought''; that is, ``rational'' means good thinking.

As an example of plain and specific language, I offer this: causal decision theory says you should take both boxes in Newcomb's problem, ``but I'd rather have a million dollars.'' \nb{The example takes a side. A few paragraphs earlier the post listed Newcomb's problem among the puzzles where the mathematics itself is disputed, and in the 2020 PhilPapers survey more philosophers chose two boxes than one.}

I suggest that you replace ``rational'' with ``foozal'' throughout the essay and see whether the connotations change.

I end: do not overuse the word. ``One receives no points merely for pronouncing it loudly.''
